\documentclass[12pt, letterpaper, english, spanish]{article}

\usepackage[T1]{fontenc}
\usepackage[utf8]{inputenc}
\usepackage[letterpaper, top = 2.5cm, bottom = 2.5cm, right = 2.5cm, left = 3cm]{geometry}
%\usepackage[myheadings]{fullpage}
\usepackage{fancyhdr}
\usepackage{lastpage}
\usepackage{enumerate}
\usepackage{enumitem}
%\usepackage{booktabs}
\usepackage{graphicx, wrapfig, setspace, booktabs}
\usepackage[font=small, labelfont=bf]{caption}
%\usepackage{bookmark}
%\usepackage[normalem,hyphens]{ulem} % Para escribir subrayado
\usepackage[normalem]{ulem} % Para escribir subrayado
\usepackage{soul} % Para escribir otros subrayados dentro de la bibliografía
\usepackage{colortbl}
\usepackage{color}
\usepackage{array}
%\usepackage[protrusion=true, expansion=true]{microtype}
\usepackage{multirow}
\usepackage{makecell}
\usepackage{tabularx}
%\usepackage{listingsutf8}
%\usepackage[usenames,dvipsnames]{color}
\usepackage[table,xcdraw]{xcolor}
\usepackage{svg}
%\usepackage{subfig}
\usepackage{subcaption}
\usepackage{pifont}
\newcommand{\cmark}{\ding{51}}%
\newcommand{\xmark}{\ding{55}}%

\hbadness=99999
\hfuzz=9999pt
\usepackage{dirtytalk}
\usepackage{authoraftertitle}
\usepackage{glossaries}
\usepackage[stable]{footmisc}

%\usepackage[noadjust, nospace, compress, sort]{cite}
%\renewcommand{\citedash}{--}
%\renewcommand{\citepunct}{,}

%Paquete para hipervinculos
\usepackage[hyphens]{url}
\usepackage[colorlinks=true]{hyperref}
\hypersetup{
    colorlinks=true,
    linkcolor=black,
    filecolor=magenta,
    urlcolor=black,
    citecolor=blue
}
%%%%%%%%%%%%%%%%%%%%%%%%%%%%%%
% Writte Cute algorithms uwu %
%%%%%%%%%%%%%%%%%%%%%%%%%%%%%%
\usepackage[linesnumbered,ruled,vlined]{algorithm2e}

%%%%%%%%%%%%%%%%%%%%%%%%%%%%%%%%%%%%%%%%%%%%
% Listings package is used to include code %
%%%%%%%%%%%%%%%%%%%%%%%%%%%%%%%%%%%%%%%%%%%%
\usepackage{caption}
\usepackage{listings}
\captionsetup[listing]{aboveskip=4pt, belowskip=2pt}
\usepackage{float}
\usepackage{csquotes}
%%%%%%%%%%%%%%%%%%%%%%%%%%%%%%%%%%%%%%%%%%%%%%%%%%%%%%%
% Listings doesn't recognize UTF-8 in files otherwise %
%%%%%%%%%%%%%%%%%%%%%%%%%%%%%%%%%%%%%%%%%%%%%%%%%%%%%%%

%%%%%%%%%%%%%%%%%%%%%%%%%%%%
% Format for included code %
%%%%%%%%%%%%%%%%%%%%%%%%%%%%
\definecolor{maroon}{cmyk}{0, 0.87, 0.68, 0.32}
\definecolor{halfgray}{gray}{0.55}
\definecolor{ipython_frame}{RGB}{207, 207, 207}
\definecolor{ipython_bg}{RGB}{247, 247, 247}
\definecolor{ipython_red}{RGB}{186, 33, 33}
\definecolor{ipython_green}{RGB}{0, 128, 0}
\definecolor{ipython_cyan}{RGB}{64, 128, 128}
\definecolor{ipython_purple}{RGB}{170, 34, 255}
\definecolor{ipython_cyan2}{RGB}{0, 150, 150}
\definecolor{ipython_orange}{RGB}{255, 93, 0}
\definecolor{ipython_magenta}{RGB}{150, 0, 150}
%%% COLORES
\definecolor{ConsoleGray}{RGB}{51, 51, 51}
\definecolor{ConsoleOrange}{RGB}{252, 192, 30}
\definecolor{ConsoleBlack}{RGB}{0,0,0}
\definecolor{ConsoleWhite}{RGB}{255,255,255}
\definecolor{MatlabBG}{rgb}{0.95,0.95,0.95}
\definecolor{codeBg}{rgb}{0,0,0}
\definecolor{mygray}{rgb}{0.85,0.85,0.85}

% Obtenidos del esquema por defecto de ansi2html (python)
\definecolor{ansi30}{HTML}{000316} % #000316 - Azul muy oscuro (casi negro)
\definecolor{ansi31}{HTML}{AA0000} % #AA0000 - Rojo oscuro
\definecolor{ansi32}{HTML}{00AA00} % #00AA00 - Verde medio
\definecolor{ansi33}{HTML}{AA5500} % #AA5500 - Naranja/marrón
\definecolor{ansi34}{HTML}{0000AA} % #0000AA - Azul oscuro
\definecolor{ansi35}{HTML}{E850A8} % #E850A8 - Rosa fuerte
\definecolor{ansi36}{HTML}{00AAAA} % #00AAAA - Cian medio
\definecolor{ansi37}{HTML}{F5F1DE} % #F5F1DE - Beige muy claro
\definecolor{ansi38}{HTML}{7F7F7F} % #7F7F7F - Gris medio
\definecolor{ansi39}{HTML}{FF0000} % #FF0000 - Rojo brillante
\definecolor{ansi310}{HTML}{00FF00} % #00FF00 - Verde brillante
\definecolor{ansi311}{HTML}{FFFF00} % #FFFF00 - Amarillo brillante
\definecolor{ansi312}{HTML}{5C5CFF} % #5C5CFF - Azul violáceo
\definecolor{ansi313}{HTML}{FF00FF} % #FF00FF - Magenta
\definecolor{ansi314}{HTML}{00FFFF} % #00FFFF - Cian brillante
\definecolor{ansi315}{HTML}{FFFFFF} % #FFFFFF - Blanco

\lstset{
    breaklines=true,
    extendedchars=true,
    stepnumber=5,
    literate=
    {á}{{\'a}}1 {é}{{\'e}}1 {í}{{\'i}}1 {ó}{{\'o}}1 {ú}{{\'u}}1
    {Á}{{\'A}}1 {É}{{\'E}}1 {Í}{{\'I}}1 {Ó}{{\'O}}1 {Ú}{{\'U}}1
    {à}{{\`a}}1 {è}{{\`e}}1 {ì}{{\`i}}1 {ò}{{\`o}}1 {ù}{{\`u}}1
    {À}{{\`A}}1 {È}{{\'E}}1 {Ì}{{\`I}}1 {Ò}{{\`O}}1 {Ù}{{\`U}}1
    {ä}{{\"a}}1 {ë}{{\"e}}1 {ï}{{\"i}}1 {ö}{{\"o}}1 {ü}{{\"u}}1
    {Ä}{{\"A}}1 {Ë}{{\"E}}1 {Ï}{{\"I}}1 {Ö}{{\"O}}1 {Ü}{{\"U}}1
    {â}{{\^a}}1 {ê}{{\^e}}1 {î}{{\^i}}1 {ô}{{\^o}}1 {û}{{\^u}}1
    {Â}{{\^A}}1 {Ê}{{\^E}}1 {Î}{{\^I}}1 {Ô}{{\^O}}1 {Û}{{\^U}}1
    {œ}{{\oe}}1 {Œ}{{\OE}}1 {æ}{{\ae}}1 {Æ}{{\AE}}1 {ß}{{\ss}}1
    {ç}{{\c c}}1 {Ç}{{\c C}}1 {ø}{{\o}}1 {å}{{\r a}}1 {Å}{{\r A}}1
    {€}{{\EUR}}1 {£}{{\pounds}}1
}
\iffalse%
\lstset{
    language=Matlab, % choose the language of the code
    basicstyle=\fontfamily{pcr}\selectfont\footnotesize\color{black},
    keywordstyle=\color[rgb]{0,0.6,0.6}\bfseries, % style for keywords
    numbers=left, % where to put the line-numbers
    backgroundcolor=\color[rgb]{0.98,0.98,0.98},
    showspaces=false, % show spaces adding particular underscores
    showstringspaces=false, % underline spaces within strings
    showtabs=false, % show tabs within strings adding particular underscores
    frame=single, % adds a frame around the code
    tabsize=4, % sets default tabsize to 2 spaces
    rulesepcolor=\color[RGB]{148, 150, 152},
    rulecolor=\color[RGB]{34, 30, 31},
    breaklines=true, % sets automatic line breaking
    breakatwhitespace=false,
    inputencoding=utf8,
    extendedchars=true,
    %commentstyle= \color[rgb]{0, 0.6, 0.6}\itshape,
    commentstyle= \color[RGB]{255, 93, 0}\itshape,
    firstnumber=1,  
    numberstyle=\tiny\color{blue}, 
    %stringstyle=\color[RGB]{148, 0, 210}\itshape,
    stringstyle=\color[rgb]{0.6, 0, 0.6},
    stepnumber=1 % los números van de 5 en 5
}
\fi
\lstdefinelanguage{matlab}{
    morekeywords={linspace,fftshift,cos,stem,xlim,ylim,xlabel,ylabel,fix,legend,access,and,break,class,continue,def,del,elif,else,except,exec,finally,for,from,global,if,import,in,is,lambda,not,or,pass,print,raise,return,try,while},
    morekeywords=[2]{abs,all,any,basestring,bin,bool,bytearray,callable,chr,classmethod,cmp,compile,complex,delattr,dict,dir,divmod,enumerate,eval,execfile,file,filter,float,format,frozenset,getattr,globals,hasattr,hash,help,hex,id,input,int,isinstance,issubclass,iter,len,list,locals,long,map,max,memoryview,min,next,object,oct,open,ord,pow,property,range,raw_input,reduce,reload,repr,reversed,round,set,setattr,slice,sorted,staticmethod,str,sum,super,tuple,type,unichr,unicode,vars,xrange,zip,apply,buffer,coerce,intern},
    sensitive=true,
    morecomment=[l]\%,
    morestring=[b]',
    morestring=[b]",
    morestring=[s]{'''}{'''},
    morestring=[s]{"""}{"""},
    morestring=[s]{r'}{'},
    morestring=[s]{r"}{"},
    morestring=[s]{r'''}{'''},
    morestring=[s]{r"""}{"""},
    morestring=[s]{u'}{'},
    morestring=[s]{u"}{"},
    morestring=[s]{u'''}{'''},
    morestring=[s]{u"""}{"""},
    % {replace}{replacement}{lenght of replace}
    % *{-}{-}{1} will not replace in comments and so on
    literate=
    {á}{{\'a}}1 {é}{{\'e}}1 {í}{{\'i}}1 {ó}{{\'o}}1 {ú}{{\'u}}1
    {Á}{{\'A}}1 {É}{{\'E}}1 {Í}{{\'I}}1 {Ó}{{\'O}}1 {Ú}{{\'U}}1
    {à}{{\`a}}1 {è}{{\`e}}1 {ì}{{\`i}}1 {ò}{{\`o}}1 {ù}{{\`u}}1
    {À}{{\`A}}1 {È}{{\'E}}1 {Ì}{{\`I}}1 {Ò}{{\`O}}1 {Ù}{{\`U}}1
    {ä}{{\"a}}1 {ë}{{\"e}}1 {ï}{{\"i}}1 {ö}{{\"o}}1 {ü}{{\"u}}1
    {Ä}{{\"A}}1 {Ë}{{\"E}}1 {Ï}{{\"I}}1 {Ö}{{\"O}}1 {Ü}{{\"U}}1
    {â}{{\^a}}1 {ê}{{\^e}}1 {î}{{\^i}}1 {ô}{{\^o}}1 {û}{{\^u}}1
    {Â}{{\^A}}1 {Ê}{{\^E}}1 {Î}{{\^I}}1 {Ô}{{\^O}}1 {Û}{{\^U}}1
    {œ}{{\oe}}1 {Œ}{{\OE}}1 {æ}{{\ae}}1 {Æ}{{\AE}}1 {ß}{{\ss}}1
    {ç}{{\c c}}1 {Ç}{{\c C}}1 {ø}{{\o}}1 {å}{{\r a}}1 {Å}{{\r A}}1
    {€}{{\EUR}}1 {£}{{\pounds}}1
    %
    {^}{{{\color{ipython_purple}\^{}}}}1
    {=}{{{\color{ipython_purple}=}}}1
    %
    {+}{{{\color{ipython_purple}+}}}1
    {*}{{{\color{ipython_purple}$^\ast$}}}1
    {/}{{{\color{ipython_purple}/}}}1
    %
    {+=}{{{+=}}}1
    {-=}{{{-=}}}1
    {*=}{{{$^\ast$=}}}1
    {/=}{{{/=}}}1,
    %
    literate=
    *{-}{{{\color{ipython_purple}-}}}1
    {?}{{{\color{ipython_purple}?}}}1
    {^}{{{\color{ipython_purple}\^{}}}}1
    {=}{{{\color{ipython_purple}=}}}1
    {+}{{{\color{ipython_purple}+}}}1
    {*}{{{\color{ipython_purple}$^\ast$}}}1
    {/}{{{\color{ipython_purple}/}}}1
    {+=}{{{+=}}}1
    {-=}{{{-=}}}1
    {*=}{{{$^\ast$=}}}1
    {/=}{{{/=}}}1,
    %
    identifierstyle=\color{black}\ttfamily,
    commentstyle=\color{ipython_orange}\ttfamily,
    stringstyle=\color{ipython_red}\ttfamily,
    keepspaces=true,
    showspaces=false,
    showstringspaces=false,
    rulecolor=\color{ipython_frame},
    frame=single,
    frameround={t}{t}{t}{t},
    framexleftmargin=6mm,
    numbers=left,
    numberstyle=\tiny\color{halfgray},
    backgroundcolor=\color{ipython_bg},
    % extendedchars=true,
    basicstyle=\scriptsize\ttfamily,%\footnotesize\ttfamily
    keywordstyle=\color{ipython_cyan2}\ttfamily,
}

\lstdefinelanguage{python}{
    morekeywords={access,and,break,class,continue,def,del,elif,else,except,exec,finally,for,from,global,if,import,in,is,lambda,not,or,pass,print,raise,return,try,while},
    morekeywords=[2]{abs,all,any,basestring,bin,bool,bytearray,callable,chr,classmethod,cmp,compile,complex,delattr,dict,dir,divmod,enumerate,eval,execfile,file,filter,float,format,frozenset,getattr,globals,hasattr,hash,help,hex,id,input,int,isinstance,issubclass,iter,len,list,locals,long,map,max,memoryview,min,next,object,oct,open,ord,pow,property,range,raw_input,reduce,reload,repr,reversed,round,set,setattr,slice,sorted,staticmethod,str,sum,super,tuple,type,unichr,unicode,vars,xrange,zip,apply,buffer,coerce,intern},
    sensitive=true,
    morecomment=[l]\#,
    morestring=[b]',
    morestring=[b]",
    morestring=[s]{'''}{'''},
    morestring=[s]{"""}{"""},
    morestring=[s]{r'}{'},
    morestring=[s]{r"}{"},
    morestring=[s]{r'''}{'''},
    morestring=[s]{r"""}{"""},
    morestring=[s]{u'}{'},
    morestring=[s]{u"}{"},
    morestring=[s]{u'''}{'''},
    morestring=[s]{u"""}{"""},
    % {replace}{replacement}{lenght of replace}
    % *{-}{-}{1} will not replace in comments and so on
    literate=
    {á}{{\'a}}1 {é}{{\'e}}1 {í}{{\'i}}1 {ó}{{\'o}}1 {ú}{{\'u}}1
    {Á}{{\'A}}1 {É}{{\'E}}1 {Í}{{\'I}}1 {Ó}{{\'O}}1 {Ú}{{\'U}}1
    {à}{{\`a}}1 {è}{{\`e}}1 {ì}{{\`i}}1 {ò}{{\`o}}1 {ù}{{\`u}}1
    {À}{{\`A}}1 {È}{{\'E}}1 {Ì}{{\`I}}1 {Ò}{{\`O}}1 {Ù}{{\`U}}1
    {ä}{{\"a}}1 {ë}{{\"e}}1 {ï}{{\"i}}1 {ö}{{\"o}}1 {ü}{{\"u}}1
    {Ä}{{\"A}}1 {Ë}{{\"E}}1 {Ï}{{\"I}}1 {Ö}{{\"O}}1 {Ü}{{\"U}}1
    {â}{{\^a}}1 {ê}{{\^e}}1 {î}{{\^i}}1 {ô}{{\^o}}1 {û}{{\^u}}1
    {Â}{{\^A}}1 {Ê}{{\^E}}1 {Î}{{\^I}}1 {Ô}{{\^O}}1 {Û}{{\^U}}1
    {œ}{{\oe}}1 {Œ}{{\OE}}1 {æ}{{\ae}}1 {Æ}{{\AE}}1 {ß}{{\ss}}1
    {ç}{{\c c}}1 {Ç}{{\c C}}1 {ø}{{\o}}1 {å}{{\r a}}1 {Å}{{\r A}}1
    {€}{{\EUR}}1 {£}{{\pounds}}1
    %
    {^}{{{\color{ipython_purple}\^{}}}}1
    {=}{{{\color{ipython_purple}=}}}1
    %
    {+}{{{\color{ipython_purple}+}}}1
    {*}{{{\color{ipython_purple}$^\ast$}}}1
    {/}{{{\color{ipython_purple}/}}}1
    %
    {+=}{{{+=}}}1
    {-=}{{{-=}}}1
    {*=}{{{$^\ast$=}}}1
    {/=}{{{/=}}}1,
    %
    literate=
    *{-}{{{\color{ipython_purple}-}}}1
    {?}{{{\color{ipython_purple}?}}}1
    {^}{{{\color{ipython_purple}\^{}}}}1
    {=}{{{\color{ipython_purple}=}}}1
    {+}{{{\color{ipython_purple}+}}}1
    {*}{{{\color{ipython_purple}$^\ast$}}}1
    {/}{{{\color{ipython_purple}/}}}1
    {+=}{{{+=}}}1
    {-=}{{{-=}}}1
    {*=}{{{$^\ast$=}}}1
    {/=}{{{/=}}}1,
    %
    identifierstyle=\color{black}\ttfamily,
    commentstyle=\color{ipython_cyan}\ttfamily,
    stringstyle=\color{ipython_red}\ttfamily,
    keepspaces=true,
    showspaces=false,
    showstringspaces=false,
    rulecolor=\color{ipython_frame},
    frame=single,
    frameround={t}{t}{t}{t},
    framexleftmargin=6mm,
    numbers=left,
    numberstyle=\tiny\color{halfgray},
    backgroundcolor=\color{ipython_bg},
    extendedchars=true,
    basicstyle=\scriptsize\ttfamily,%\scriptsize,
    keywordstyle=\color{ipython_green}\ttfamily,
}

%%%%%%%%%%%%%%%%%%%%%%%%
% Fancy Caption Format %
%%%%%%%%%%%%%%%%%%%%%%%%

\DeclareCaptionFont{white}{ \color{white} }
\DeclareCaptionFormat{listing}{
  \colorbox[cmyk]{0.9, 0.35, 0.35, 0.01 }{
    \parbox{1\linewidth}{\hspace{0pt}#1#2#3}
  }\vskip+1pt
}
% Not used in this proyect
% \captionsetup[lstlisting]{format=listing, labelfont=white, textfont=white, singlelinecheck=false, margin=15pt, font={bf,footnotesize}}

% CMD (DOS)
\lstdefinestyle{DOS} {
    backgroundcolor=\color{black},
    basicstyle=\footnotesize\color{white}\ttfamily
}
% Consola Putty usada por GNS3
\lstdefinestyle{SolarPutty} {
    backgroundcolor=\color{ConsoleGray},
    basicstyle=\footnotesize\color{ConsoleOrange}\ttfamily
}
% Consola Putty usada por GNS3 (más chico)
\lstdefinestyle{SolarPutty2} {
    backgroundcolor=\color{ConsoleGray},
    basicstyle=\scriptsize\color{ConsoleOrange}\ttfamily,
    columns=flexible,
    %breaklines=true,
    %breakatwhitespace=true,
    %postbreak=\mbox{\textcolor{ConsoleOrange}{$\hookrightarrow$}\space}
}

%%% MOSTRAR CODIGO (Usando minted)
\usepackage[newfloat,outputdir=out]{minted}
% newfloat para que cuente como lstlisting en la tabla de contenidos
% outputdir=out para que compile correctamente
% IMPORTANTE!! Invocar LaTeX con -shell-escape
\usemintedstyle[python]{monokai}
\usemintedstyle[c]{manni}
\setminted[python] {
    breaklines=true,
    encoding=utf8,
    fontsize=\footnotesize,
    bgcolor=codeBg
}
\setminted[c] {
    breaklines=true,
    encoding=utf8,
    fontsize=\footnotesize
}
\setminted[matlab] {
    bgcolor=MatlabBG,
    linenos=true,
    fontsize=\footnotesize
}


%%% MOSTRAR CODIGO (Con Verbatim y colorbox)
\usepackage[most]{tcolorbox}
\usepackage{verbatim}
\usepackage{fancyvrb}
\fvset{commandchars=\\\{\}} % Permite escapar códigos de LaTex
% Mostrar código como en la consola con coloreado ANSI
\newtcolorbox{customText}[1][]{
  colback=ConsoleBlack,        % Color de fondo
  coltext=ConsoleWhite,        % Color del texto
  boxrule=0pt,         % Sin borde
  sharp corners,       % Esquinas rectas
  breakable,           % Permite romper el bloque en varias páginas si es largo
  left=2mm, right=2mm, top=2mm, bottom=2mm, % Márgenes internos
  fontupper=\ttfamily\scriptsize % Fuente y tamaño del texto
}
\newtcolorbox{customText2}[1][]{
  colback=ConsoleGray,        % Color de fondo
  coltext=ConsoleOrange,        % Color del texto
  boxrule=0pt,         % Sin borde
  sharp corners,       % Esquinas rectas
  breakable,           % Permite romper el bloque en varias páginas si es largo
  left=0mm, right=0mm, top=0mm, bottom=0mm, % Márgenes internos
  fontupper=\ttfamily\footnotesize % Fuente y tamaño del texto
}
\newtcolorbox{customText3}[1][]{
  colback=ConsoleBlack,        % Color de fondo
  coltext=ConsoleWhite,        % Color del texto
  boxrule=0pt,         % Sin borde
  sharp corners,       % Esquinas rectas
  breakable,           % Permite romper el bloque en varias páginas si es largo
  left=2mm, right=2mm, top=2mm, bottom=2mm, % Márgenes internos
  fontupper=\ttfamily\tiny % Fuente y tamaño del texto
}
% Código en linea con fondo gris en un cuadrado con bordes redondos
\newtcbox{\lsti}[1][]{
    on line,
    box align=base,
    enhanced,
    colback=mygray,
    colframe=mygray,
    boxrule=0pt,
    arc=2mm,
    left=1pt,
    right=1pt,
    top=-1pt,
    bottom=-1pt,
    fontupper=\ttfamily,
    nobeforeafter,
    #1
}


%Para qué los subtítulos aparezcan en español
\usepackage[spanish,es-tabla,es-noshorthands,es-nolists,es-lcroman]{babel}
%\usepackage[spanish]{babel}
%\usepackage[utf8]{inputenc}
%\usepackage{sectsty}
\usepackage{titling, lipsum}
\usepackage{tabularx}
\usepackage{float}
\usepackage{listings}
%\usepackage[numbers, sort&compress]{natbib}
\usepackage{anyfontsize}
\usepackage{parskip}
%\usepackage{setspace}

%Herramientas matemáticas
\usepackage{amsmath}
\usepackage{breqn}
\usepackage{siunitx}
\usepackage{amssymb}
\usepackage{mathtools}
\usepackage{mathrsfs}
\usepackage{textcomp}
\usepackage{gensymb}

%%%%%%%%%%%%%%%%%%%%%%%%%%%%%%
% REFERENCIAS BIBLIOGRÁFICAS %
%%%%%%%%%%%%%%%%%%%%%%%%%%%%%%
\usepackage[backend=biber,datamodel=citasModeloUSM,style=ieee,citestyle=numeric-comp,natbib=true,minnames=1,maxnames=3,alldates=comp,alldatesusetime=true,alltimes=24h,dateera=christian,dateeraauto=1000,timezones=true,seconds=true,block=none,sorting=none]{biblatex}
%\usepackage[backend=biber,style=ieee]{biblatex}
\usepackage{citasEstiloUSM} % Estilo de citas de la USM
\usepackage{xurl}
%\renewcommand{\thesubsection}{\thesection.\alph{subsection}}

\newenvironment{poliabstract}[1]
  {%
    \renewcommand{\abstractname}{#1}% Cambia el nombre del abstract
    \begin{abstract}%
      %\noindent% Elimina sangría inicial
      \setlength{\leftskip}{0pt}% Asegura que el margen izquierdo sea cero
      \setlength{\rightskip}{0pt}% Asegura que el margen derecho sea cero
  }
  {\end{abstract}}


\newcommand{\HRule}[1]{\rule{\linewidth}{#1}}
\onehalfspacing
\setcounter{tocdepth}{5}
\setcounter{secnumdepth}{5}

% --- NUMERACIÓN: incluir número de la parte antes de las secciones ---
\usepackage{chngcntr}                % para resetear contadores con facilidad

% forzar que \thepart sea arábigo (opcional, evita I, II en vez de 1,2)
\renewcommand{\thepart}{\arabic{part}}

% que las secciones se reinicien al cambiar de parte
\counterwithin{section}{part}

% numeración visible: Parte.Sección.Subsección...
\renewcommand{\thesection}{\thepart.\arabic{section}}
\renewcommand{\thesubsection}{\thesection.\arabic{subsection}}
\renewcommand{\thesubsubsection}{\thesubsection.\arabic{subsubsection}}

% actualizar numeración de figuras/tablas/ecuaciones para que usen la nueva sección
\renewcommand{\thefigure}{\thesection.\arabic{figure}}
\renewcommand{\thetable}{\thesection.\arabic{table}}
\renewcommand{\theequation}{\thesection.\arabic{equation}}

% opcional: hacer que figure/table/equation se reinicien con la sección
\counterwithin{figure}{section}
\counterwithin{table}{section}
\counterwithin{equation}{section}

% agregar partes al índice de contenidos con numeración
\newcommand{\parttoc}[1]{%
  \part*{#1}%
  \refstepcounter{part}%
  \phantomsection%
  \addcontentsline{toc}{part}{\protect\numberline{\thepart}#1}%
}

% ---------- AJUSTES DE TOC (evitar solapamiento de números largos) ----------
\usepackage{tocloft}

% Ajusta anchuras de la “columna de números” en el TOC (tunea estos valores)
\setlength{\cftpartnumwidth}{1.5em}     % ancho columna para "Partes" (número)
\setlength{\cftsecnumwidth}{2em}      % ancho columna para las secciones
\setlength{\cftsubsecnumwidth}{2.9em}   % ancho columna para las subsecciones
\setlength{\cftsubsubsecnumwidth}{3.8em}% ancho para subsubsecciones si las listaste
\setlength{\cftfignumwidth}{2.9em}      % ancho columna para números de figuras
\setlength{\cfttabnumwidth}{2.9em}      % ancho columna para números de tablas

% Si deseas, ajusta el espacio entre número y texto
%\setlength{\cftbeforesecskip}{0pt}      % espacio antes de cada entrada de sección
% --------------------------------------------------------------------------


\providecommand{\keywordses}[1]
{
  \small	
  \textbf{\textit{Palabras Clave: }} #1
}

\providecommand{\keywordsen}[1]
{
  \small	
  \textbf{\textit{Keywords: }} #1
}

\usepackage{datetime}
\newdateformat{monthyeardate}{%
\monthname[\THEMONTH] \THEYEAR}

\newdateformat{monthsyeardate}{%
\monthname[\THEMONTH]-\THEYEAR}

  

\newdateformat{mifecha}{Valparaíso, a \THEDAY~de \monthname[\THEMONTH] de \THEYEAR}
\makeatletter
\renewcommand{\monthnamespanish}[1][\month]{%
  \@orgargctr=#1\relax
  \ifcase\@orgargctr
    \PackageError{datetime}{Invalid Month number \the\@orgargctr}{%
      Month numbers should go from 1 to 12}%
    \or Enero%
    \or Febrero%
    \or Marzo%
    \or Abril%
    \or Mayo%
    \or Junio%
    \or Julio%
    \or Agosto%
    \or Septiembre%
    \or Octubre%
    \or Noviembre%
    \or Diciembre%
    \else \PackageError{datetime}{Invalid Month number \the\@orgargctr}{%
      Month numbers should go from 1 to 12}%
    \fi%
}%
\makeatother



\fancyfoot[C]{}
\fancyfoot[R]{\thepage}
\renewcommand{\headrulewidth}{0pt}
\renewcommand{\footrulewidth}{0pt}



\newif\ifmargins
%\marginstrue % COMENTAR CUANDO SE RELLENE EL DOCUMENTO

\ifmargins
\geometry{showframe}
\usepackage{fgruler}
\fi

\newif\ifbib
%\ifbibtrue

\ifbib
\renewcommand{\bibname}{Bibliografía}
\renewcommand{\refname}{Bibliografía}
\def\refname{Bibliografía}
\renewcommand{\bibsection}{\section*{Bibliografía}}
\fi

\headheight = 1.5cm
\textheight = 22cm
%\headsep = 0.5cm
%\topmargin = 1.5cm

%\renewcommand{\baselinestretch}{1.5} 
\linespread{1.5}
\newlength\tindent
\setlength{\tindent}{\parindent}
\setlength{\parindent}{0pt}
\renewcommand{\indent}{\hspace*{\tindent}}

% Environment renaming
\addto\captionsspanish{
    \renewcommand{\contentsname}%
        {Índice de contenidos}%
}
%\renewcommand{\listtablename}{Índice de tablas}
%\renewcommand{\tablename}{Tabla}
\makeatletter
\AtBeginDocument{%
\renewcommand\lstlistlistingname{Índice de códigos y algoritmos}
 \let\c@algocf\c@lstlisting
}
\renewcommand{\algocf@list}{lol}%
\renewcommand*\l@algocf{\@dottedtocline{1}{1.5em}{2.3em}}
\makeatother

\renewcommand{\listingname}{Código}
\renewcommand{\lstlistingname}{Código}
\newcommand*\lstinputpath[1]{\lstset{inputpath=#1}}

\newcommand\blfootnote[1]{%
    \begingroup
    \renewcommand\thefootnote{}\footnote{#1}%
    \addtocounter{footnote}{-1}%
    \endgroup
}

% Figures / Tables / Equations incluirán parte.sección.XXX
\renewcommand{\thefigure}{\thesection.\arabic{figure}}
\renewcommand{\thetable}{\thesection.\arabic{table}}
\renewcommand{\theequation}{\thesection.\arabic{equation}}
%\renewcommand{\thelstlisting}{\thesection.\arabic{lstlisting}}
%\makeatletter
%\renewcommand{\thealgocf}{\thesection.\arabic{algocf}}
%\makeatother


% Table colours
\definecolor{LightGrey}{gray}{0.9}
\definecolor{LightCyan}{rgb}{0.88,1,1}
\definecolor{White}{gray}{0}
\definecolor{colour1}{HTML}{7B84FC}
\definecolor{colour2}{HTML}{7AD6FD}
\definecolor{colour3}{HTML}{79FBD6}
\definecolor{colour4}{HTML}{78F87F}
\definecolor{colour5}{HTML}{D5FA80}
\definecolor{colour6}{HTML}{FFFA81}
\definecolor{colour7}{HTML}{FED37F}
\definecolor{colour8}{HTML}{FD7F7C}

% Path config
\graphicspath{{figures/}{logos/}}
\lstinputpath{code/}

\makeatletter
\def\include@path{{source/}}
\def\input@path{{source/}}
\makeatother


% New environments
\newenvironment{dedicatory}{
\clearpage
\null
%\vspace{\fill}
\vspace{\stretch{1}}
\begin{flushright}\itshape}{
\end{flushright}
\vspace{\stretch{.1}}\null
}


\let\thanks\undefined
\newenvironment{thanks}{
	\vspace*{\fill}
	\if@en
	    \begin{center}
		    \section*{Acknowledgments}
		\end{center}
	\else
	    \begin{center}
		    \centering\section*{Agradecimientos}
		\end{center}
	\fi
	\vspace*{\fill}
}
{\vfill\null}


\let\thanksnf\undefined
\newenvironment{thanksnf}{
	%\vspace*{\fill}
	\if@en
	    \begin{center}
		    \section*{Acknowledgments}
		\end{center}
	\else
	    \begin{center}
		    \centering\section*{Agradecimientos}
		\end{center}
	\fi
	%\vspace*{\fill}
}
{\vfill\null}


\makeatletter
\newcommand\frontmatter{%
    \cleardoublepage
  %\@mainmatterfalse
  \pagenumbering{roman}}

\newcommand\mainmatter{%
    \cleardoublepage
 % \@mainmattertrue
  \pagenumbering{arabic}}

%\newcommand\backmatter{%
%  \if@openright
%    \cleardoublepage
%  \else
%    \clearpage
%  \fi
%  \@mainmatterfalse
%   }
\makeatother

\newcommand*{\makecover}{%
    \begin{center}
    {\fontsize{18}{18}\selectfont \textbf{\institution}}
    \vspace{.7cm}

    {\fontsize{16}{16}\selectfont \textbf{\department}}
    \vspace{.7cm}

    {\fontsize{14}{14}\selectfont \textbf{\location}}
    \vspace{.7cm}

    \includegraphics[height=33mm]{\logo}\\
	\vspace{.7cm}

	{\fontsize{20}{20}\selectfont \textbf{\say{\subject}}}\blfootnote{\disclaimer}
	\vspace{0.7cm}

	{\fontsize{14}{14}\selectfont \textbf{\authorName}}
	\vspace{0.7cm}

	{\fontsize{12}{12}\selectfont \textbf{\details}}
    \vspace{0.7cm}

	{\fontsize{12}{12}\selectfont \textbf{\advisor}}
    \vspace{.7cm}

	{\fontsize{14}{14}\selectfont \textsc{\textbf{\monthsyeardate\examdate}}}
    \end{center}
}


\renewcommand{\glstextformat}[1]{\textit{#1}}



%%%%%%%%%%%%%%%%%%%%%%%%%%%%%%%%%%%%%%%%%%%%%
% FIGURAS DE DISPOSICIÓN Y TAMAÑO AJUSTABLE %
%%%%%%%%%%%%%%%%%%%%%%%%%%%%%%%%%%%%%%%%%%%%%
\usepackage{xfp}      % cálculos en punto flotante
\usepackage{calc}
\usepackage{xparse}
\usepackage{etoolbox} % utilidades CSV (opcional)
\makeatletter

% ---------------- \TwoImgs (simple, para 2 imágenes con misma altura calculada) ----------------
\newsavebox{\TwoImgs@BoxA}
\newsavebox{\TwoImgs@BoxB}
\newlength{\TwoImgs@maxwidth}
\newlength{\TwoImgs@sep}
\newlength{\TwoImgs@Awd}
\newlength{\TwoImgs@Aht}
\newlength{\TwoImgs@Bwd}
\newlength{\TwoImgs@Bht}
\newlength{\TwoImgs@computedheight}
\setlength{\TwoImgs@maxwidth}{0.9\textwidth}
\setlength{\TwoImgs@sep}{1em}
\newcommand{\TwoImgsSetDefaults}[2]{%
  \setlength{\TwoImgs@maxwidth}{#1}%
  \setlength{\TwoImgs@sep}{#2}%
}
\NewDocumentCommand{\TwoImgs}{ O{\TwoImgs@maxwidth} O{\TwoImgs@sep} m m m m O{} m m }{%
  \setlength{\TwoImgs@maxwidth}{#1}%
  \setlength{\TwoImgs@sep}{#2}%
  % medir imágenes
  \sbox{\TwoImgs@BoxA}{\includegraphics{#3}}%
  \sbox{\TwoImgs@BoxB}{\includegraphics{#5}}%
  \setlength{\TwoImgs@Awd}{\wd\TwoImgs@BoxA}%
  \setlength{\TwoImgs@Aht}{\ht\TwoImgs@BoxA}%
  \setlength{\TwoImgs@Bwd}{\wd\TwoImgs@BoxB}%
  \setlength{\TwoImgs@Bht}{\ht\TwoImgs@BoxB}%
  % altura común (fórmula estable)
  \setlength{\TwoImgs@computedheight}{%
    \fpeval{%
      (\strip@pt\TwoImgs@maxwidth - \strip@pt\TwoImgs@sep) /
      ( (\strip@pt\TwoImgs@Awd)/(\strip@pt\TwoImgs@Aht) + (\strip@pt\TwoImgs@Bwd)/(\strip@pt\TwoImgs@Bht) )
    }pt%
  }%
  % salida
  \begin{figure}[htbp]
    \centering
    \begin{subfigure}[b]{\dimexpr 0.5\TwoImgs@maxwidth\relax}
      \centering
      \includegraphics[height=\TwoImgs@computedheight,keepaspectratio]{#3}%
      \caption{#4}%
    \end{subfigure}%
    \hspace{\TwoImgs@sep}%
    \begin{subfigure}[b]{\dimexpr 0.5\TwoImgs@maxwidth\relax}
      \centering
      \includegraphics[height=\TwoImgs@computedheight,keepaspectratio]{#5}%
      \caption{#6}%
    \end{subfigure}
    % aquí elegimos entre \caption{long} o \caption[short]{long}
    \ifstrempty{#7}{%
      \caption{#8}%
    }{%
      \caption[#7]{#8}%
    }%
    \label{#9}
  \end{figure}
}

% ---------------- \EqualColGrid (divide el ancho en columnas iguales) ----------------
\newlength{\EqualColGrid@cellwidth}
\NewDocumentCommand{\EqualColGrid}{ O{0.9\textwidth} O{1em} m m m m m }{%
  % #1 maxwidth, #2 sep, #3 cols, #4 files CSV, #5 captions CSV, #6 caption general, #7 label
  \begingroup
  \setlength{\@tempdima}{#1}%
  \setlength{\@tempdimb}{#2}%
  \edef\ECG@cols{#3}%
  \setlength{\EqualColGrid@cellwidth}{\dimexpr (\@tempdima - (\ECG@cols-1)\@tempdimb) / \ECG@cols\relax}%
  \begin{figure}[htbp]\centering
    \newcounter{ECG@i}\setcounter{ECG@i}{0}%
    % helper to fetch nth caption
    \newcommand\EqualColGrid@getcap[1]{%
      \edef\EqualColGrid@ans{}%
      \newcounter{ECG@tmp}\setcounter{ECG@tmp}{0}%
      \renewcommand{\do##1}{\stepcounter{ECG@tmp}\ifnum\value{ECG@tmp}=#1 \edef\EqualColGrid@ans{##1}\fi}%
      \docsvlist{#5}%
      \EqualColGrid@ans
    }%
    \renewcommand{\do}[1]{%
      \stepcounter{ECG@i}%
      \edef\ECG@file{##1}%
      \edef\ECG@cap{\EqualColGrid@getcap{\theECG@i}}%
      \begin{subfigure}[b]{\EqualColGrid@cellwidth}\centering
        \includegraphics[width=\linewidth,keepaspectratio]{\ECG@file}%
        \caption{\ECG@cap}%
      \end{subfigure}%
      % separación y salto de línea al final de cada fila
      \ifnum\value{ECG@i}<\numexpr\docsvlistlen{#4}\relax
        \ifnum\value{ECG@i}<\ECG@cols
          \hspace{#2}%
        \fi
        \ifnum\value{ECG@i}=\numexpr\ECG@cols * (\numexpr(\value{ECG@i}-1)/\ECG@cols\relax)\relax
          \\[0.5em]%
        \fi
      \fi
    }%
    \docsvlist{#4}%
    \caption{#6}\label{#7}%
  \end{figure}
  \endgroup
}

% ---------------- \UniformGrid (misma altura para todas las imágenes; filas no exceden maxwidth) ----------------
\newcounter{UG@count}
\newcounter{UG@cols}
\newcounter{UG@rows}
\newlength{\UG@maxwidth}
\newlength{\UG@seph}
\newlength{\UG@cellwidth}
\newlength{\UG@H}
\NewDocumentCommand{\UniformGrid}{ O{0.9\textwidth} O{1em} m m m m m }{%
  % #1 = maxwidth, #2 = sep, #3 = cols, #4 = files CSV, #5 = captions CSV, #6 caption, #7 label
  \begingroup
  \setlength{\UG@maxwidth}{#1}%
  \setlength{\UG@seph}{#2}%
  \setcounter{UG@cols}{#3}%
  \setcounter{UG@count}{0}%
  % 1) medir y almacenar coef = w/h para cada imagen
  \@for\UG@file:=#4\do{%
    \stepcounter{UG@count}%
    % medir
    \sbox0{\includegraphics{\UG@file}}%
    \setlength{\@tempdima}{\wd0}%
    \setlength{\@tempdimb}{\ht0}%
    \expandafter\xdef\csname UG@file@\theUG@count\endcsname{\UG@file}%
    \expandafter\xdef\csname UG@w@\theUG@count\endcsname{\strip@pt\@tempdima}%
    \expandafter\xdef\csname UG@h@\theUG@count\endcsname{\strip@pt\@tempdimb}%
    \expandafter\xdef\csname UG@coef@\theUG@count\endcsname{\fpeval{(\strip@pt\@tempdima)/(\strip@pt\@tempdimb)}}%
  }%
  \edef\UG@N{\theUG@count}%
  \setcounter{UG@rows}{\numexpr (\UG@N + \value{UG@cols} -1) / \value{UG@cols} \relax}%
  % 2) para cada fila computar sumcoef y H_row, quedarnos con min(H_row)
  \edef\UG@minH{1000000}%
  \count@=1
  \loop
    \ifnum\count@>\value{UG@rows}\relax
      \breakloop
    \fi
    \edef\UG@start{\the\numexpr (\count@-1)*\value{UG@cols} + 1\relax}%
    \edef\UG@end{\the\numexpr \count@*\value{UG@cols}\relax}%
    % corregir si pasa de N
    \ifnum\UG@end>\UG@N \edef\UG@end{\UG@N}\fi
    % sumar coeficientes en la fila
    \edef\UG@sumcoeff{0}%
    \count@i=\UG@start
    \loop
      \ifnum\count@i>\UG@end \breakloop \fi
      \edef\UG@coefval{\csname UG@coef@\the\count@i\endcsname}%
      \edef\UG@sumcoeff{\fpeval{\UG@sumcoeff + \UG@coefval}}%
      \advance\count@i by 1
    \repeat
    \edef\UG@seppt{\strip@pt\UG@seph}%
    \edef\UG@maxwpt{\strip@pt\UG@maxwidth}%
    \edef\UG@Hrow{\fpeval{(\UG@maxwpt - (\UG@cols-1)*\UG@seppt) / \UG@sumcoeff}}%
    \edef\UG@minH{\fpeval{min(\UG@minH,\UG@Hrow)}}%
    \advance\count@ by 1
  \repeat
  % set length for height
  \setlength{\UG@H}{\UG@minH pt}%
  % cell width for layout (equal columns)
  \setlength{\UG@cellwidth}{\dimexpr (\UG@maxwidth - (\value{UG@cols}-1)\UG@seph) / \value{UG@cols}\relax}%
  % 3) render
  \begin{figure}[htbp]\centering
    \newcounter{UG@renderi}\setcounter{UG@renderi}{0}%
    \newcounter{UG@rendercol}\setcounter{UG@rendercol}{0}%
    % helper para nth caption
    \newcommand\UG@getcap[1]{%
      \edef\UG@ans{}%
      \newcounter{UG@tmp}\setcounter{UG@tmp}{0}%
      \renewcommand{\do##1}{\stepcounter{UG@tmp}\ifnum\value{UG@tmp}=#1 \edef\UG@ans{##1}\fi}%
      \docsvlist{#5}%
      \UG@ans
    }%
    \@for\UG@file:=#4\do{%
      \stepcounter{UG@renderi}\stepcounter{UG@rendercol}%
      \edef\UG@thisfile{\UG@file}%
      \edef\UG@thiscap{\UG@getcap{\theUG@renderi}}%
      \begin{subfigure}[b]{\UG@cellwidth}\centering
        \includegraphics[height=\UG@H,keepaspectratio]{\UG@thisfile}%
        \caption{\UG@thiscap}%
      \end{subfigure}%
      \ifnum\value{UG@renderi}<\UG@N
        \ifnum\value{UG@rendercol}<\value{UG@cols}
          \hspace{\UG@seph}%
        \else
          \\[0.6em]\setcounter{UG@rendercol}{0}%
        \fi
      \fi
    }%
    \caption{#6}\label{#7}%
  \end{figure}
  \endgroup
}
\makeatother

%% TODO ? Los métodos \EqualColGrid y \UniformGrid no funcan. No se porqué. Arreglar si quieres, no se xdddd.